\documentclass[10pt]{extarticle}
\usepackage[a4paper,margin=1.5cm]{geometry}
\usepackage[T1]{fontenc}
\usepackage{amsmath,amssymb}
\usepackage{graphicx}
\usepackage[export]{adjustbox}
\usepackage{multicol}
\usepackage{enumitem}
\usepackage{booktabs,array,tabularx}
\setlength{\columnsep}{0.7cm}
\setlength{\parindent}{0pt}
\setlength{\parskip}{3pt}
\pagestyle{empty}
\makeatletter
\newcommand{\qfigdef}[4]{\@namedef{qf@file@#1}{#2}\@namedef{qf@w@#1}{#3}\@namedef{qf@wide@#1}{#4}}
\newcommand{\qf@check}[1]{\@ifundefined{qf@file@#1}{\PackageError{score2}{Unknown figure id #1}{}}{}}
\newcommand{\qf@img}[4][]{\includegraphics[width=#2\dimexpr\@nameuse{qf@w@#4}pt\relax,max width=#3,max height=0.45\textheight,#1]{\@nameuse{qf@file@#4}}}
\newif\ifqf@full
\newif\ifqf@boxed
\newsavebox{\qf@box}
\newcommand{\qf@open}{\global\qf@fulltrue\par\global\qf@boxedtrue
  \setbox\qf@box\vbox\bgroup\hsize\textwidth\linewidth\textwidth}
\newcommand{\qf@emit}[1]{\par\egroup\global\qf@boxedfalse
  \ifdim\dimexpr\ht\qf@box+\dp\qf@box\relax<#1\textheight\noindent\usebox{\qf@box}\par\else\unvbox\qf@box\fi}
\newcommand{\qfullsplit}{\ifqf@boxed\qf@emit{0.85}\fi}
\newcommand{\qf@close}{\ifqf@boxed\qf@emit{0.45}\fi\global\qf@fullfalse}
\newcommand{\qfullbegin}{\end{multicols}\qf@open}
\newcommand{\qfullend}{\qf@close\begin{multicols}{2}}
\newcommand{\qfullstart}{\qf@open}
\newcommand{\qfulldone}{\qf@close}
\newcommand{\qfig}[2][1]{\qf@check{#2}%
  \ifqf@full\qf@blk{#1}{#2}\else
  \ifnum\@nameuse{qf@wide@#2}=1 \expandafter\@firstoftwo\else\expandafter\@secondoftwo\fi
  {\end{multicols}\begin{center}\qf@img{#1}{\textwidth}{#2}\end{center}\begin{multicols}{2}}%
  {\qf@blk{#1}{#2}}\fi}
\newcommand{\qf@blk}[2]{\par\smallskip{\centering\qf@img{#1}{\linewidth}{#2}\par}\smallskip}
\newcommand{\qfigi}[2][1]{\qf@check{#2}\qf@img[valign=c]{#1}{\linewidth}{#2}}
\makeatother
\newcommand{\papersec}[1]{\par\bigskip{\centering\Large\bfseries PAPER-#1\par}\medskip}
\newcommand{\partsec}[2]{\par\medskip\noindent{\large\bfseries PART-#1: #2}\par\smallskip}
\newcommand{\sect}[1]{\par\smallskip\noindent\textbf{#1}\par}
\newcommand{\qn}[1]{\par\smallskip\noindent\textbf{#1.}~}
\newcommand{\paraq}[2]{\par\smallskip\noindent\textbf{Paragraph for Questions #1 and #2}\par}
\newenvironment{qoptions}{\begin{itemize}[leftmargin=1.8em,itemsep=1pt,topsep=2pt,label={}]}{\end{itemize}}
\newcommand{\opt}[2]{\item[(#1)] #2}
% figure registry (generated by score2_tools.py finalize)
\qfigdef{p1_phy_s1_q01_1}{JEE Advanced (2021 Series) - Test 13 (31 May 2021)-img/p1_phy_s1_q01_1.png}{158.2}{0}
\qfigdef{p1_phy_s1_q02_1}{JEE Advanced (2021 Series) - Test 13 (31 May 2021)-img/p1_phy_s1_q02_1.png}{131.0}{0}
\qfigdef{p1_phy_s1_q03_1}{JEE Advanced (2021 Series) - Test 13 (31 May 2021)-img/p1_phy_s1_q03_1.png}{108.5}{0}
\qfigdef{p1_phy_s1_q04_1}{JEE Advanced (2021 Series) - Test 13 (31 May 2021)-img/p1_phy_s1_q04_1.png}{70.1}{0}
\qfigdef{p1_phy_s1_q05_1}{JEE Advanced (2021 Series) - Test 13 (31 May 2021)-img/p1_phy_s1_q05_1.png}{74.4}{0}
\qfigdef{p1_phy_s1_q06_1}{JEE Advanced (2021 Series) - Test 13 (31 May 2021)-img/p1_phy_s1_q06_1.png}{59.9}{0}
\qfigdef{p1_phy_s1_q07_1}{JEE Advanced (2021 Series) - Test 13 (31 May 2021)-img/p1_phy_s1_q07_1.png}{103.6}{0}
\qfigdef{p1_phy_s2_pre_1}{JEE Advanced (2021 Series) - Test 13 (31 May 2021)-img/p1_phy_s2_pre_1.png}{50.1}{0}
\qfigdef{p1_phy_s2_pre_2}{JEE Advanced (2021 Series) - Test 13 (31 May 2021)-img/p1_phy_s2_pre_2.png}{30.4}{0}
\qfigdef{p1_phy_s2_pre_3}{JEE Advanced (2021 Series) - Test 13 (31 May 2021)-img/p1_phy_s2_pre_3.png}{26.5}{0}
\qfigdef{p1_phy_s2_pre_4}{JEE Advanced (2021 Series) - Test 13 (31 May 2021)-img/p1_phy_s2_pre_4.png}{49.0}{0}
\qfigdef{p1_phy_s2_q03_1}{JEE Advanced (2021 Series) - Test 13 (31 May 2021)-img/p1_phy_s2_q03_1.png}{117.5}{0}
\qfigdef{p1_phy_s2_q03_2}{JEE Advanced (2021 Series) - Test 13 (31 May 2021)-img/p1_phy_s2_q03_2.png}{116.1}{0}
\qfigdef{p1_phy_s2_q03_3}{JEE Advanced (2021 Series) - Test 13 (31 May 2021)-img/p1_phy_s2_q03_3.png}{93.9}{0}
\qfigdef{p1_phy_s2_q03_4}{JEE Advanced (2021 Series) - Test 13 (31 May 2021)-img/p1_phy_s2_q03_4.png}{113.6}{0}
\qfigdef{p1_phy_s3_q03_1}{JEE Advanced (2021 Series) - Test 13 (31 May 2021)-img/p1_phy_s3_q03_1.png}{107.6}{0}
\qfigdef{p1_phy_s3_q04_1}{JEE Advanced (2021 Series) - Test 13 (31 May 2021)-img/p1_phy_s3_q04_1.png}{156.2}{0}
\qfigdef{p1_chem_s1_q03_1}{JEE Advanced (2021 Series) - Test 13 (31 May 2021)-img/p1_chem_s1_q03_1.png}{309.4}{1}
\qfigdef{p1_chem_s1_q07_A}{JEE Advanced (2021 Series) - Test 13 (31 May 2021)-img/p1_chem_s1_q07_A.png}{43.1}{0}
\qfigdef{p1_chem_s1_q07_B}{JEE Advanced (2021 Series) - Test 13 (31 May 2021)-img/p1_chem_s1_q07_B.png}{46.5}{0}
\qfigdef{p1_chem_s1_q07_C}{JEE Advanced (2021 Series) - Test 13 (31 May 2021)-img/p1_chem_s1_q07_C.png}{45.3}{0}
\qfigdef{p1_chem_s1_q07_D}{JEE Advanced (2021 Series) - Test 13 (31 May 2021)-img/p1_chem_s1_q07_D.png}{45.8}{0}
\qfigdef{p1_chem_s3_q03_1}{JEE Advanced (2021 Series) - Test 13 (31 May 2021)-img/p1_chem_s3_q03_1.png}{90.8}{0}
\qfigdef{p2_phy_s1_q01_1}{JEE Advanced (2021 Series) - Test 13 (31 May 2021)-img/p2_phy_s1_q01_1.png}{79.6}{0}
\qfigdef{p2_phy_s1_q03_B}{JEE Advanced (2021 Series) - Test 13 (31 May 2021)-img/p2_phy_s1_q03_B.png}{145.6}{0}
\qfigdef{p2_phy_s2_q02_1}{JEE Advanced (2021 Series) - Test 13 (31 May 2021)-img/p2_phy_s2_q02_1.png}{145.0}{0}
\qfigdef{p2_phy_s2_q03_1}{JEE Advanced (2021 Series) - Test 13 (31 May 2021)-img/p2_phy_s2_q03_1.png}{104.6}{0}
\qfigdef{p2_phy_s2_q04_1}{JEE Advanced (2021 Series) - Test 13 (31 May 2021)-img/p2_phy_s2_q04_1.png}{100.7}{0}
\qfigdef{p2_phy_s2_q05_1}{JEE Advanced (2021 Series) - Test 13 (31 May 2021)-img/p2_phy_s2_q05_1.png}{130.5}{0}
\qfigdef{p2_phy_s3_para1_2_1}{JEE Advanced (2021 Series) - Test 13 (31 May 2021)-img/p2_phy_s3_para1_2_1.png}{60.1}{0}
\qfigdef{p2_phy_s4_q02_1}{JEE Advanced (2021 Series) - Test 13 (31 May 2021)-img/p2_phy_s4_q02_1.png}{100.2}{0}
\qfigdef{p2_phy_s5_q01_1}{JEE Advanced (2021 Series) - Test 13 (31 May 2021)-img/p2_phy_s5_q01_1.png}{65.9}{0}
\qfigdef{p2_phy_s5_q02_1}{JEE Advanced (2021 Series) - Test 13 (31 May 2021)-img/p2_phy_s5_q02_1.png}{181.3}{0}
\qfigdef{p2_phy_s5_q03_1}{JEE Advanced (2021 Series) - Test 13 (31 May 2021)-img/p2_phy_s5_q03_1.png}{128.0}{0}
\qfigdef{p2_phy_s5_q04_1}{JEE Advanced (2021 Series) - Test 13 (31 May 2021)-img/p2_phy_s5_q04_1.png}{60.9}{0}
\qfigdef{p2_phy_s5_q05_1}{JEE Advanced (2021 Series) - Test 13 (31 May 2021)-img/p2_phy_s5_q05_1.png}{105.3}{0}
\qfigdef{p2_phy_s5_q06_1}{JEE Advanced (2021 Series) - Test 13 (31 May 2021)-img/p2_phy_s5_q06_1.png}{48.5}{0}
\qfigdef{p2_phy_s5_q07_1}{JEE Advanced (2021 Series) - Test 13 (31 May 2021)-img/p2_phy_s5_q07_1.png}{102.1}{0}
\qfigdef{p2_phy_s5_q10_1}{JEE Advanced (2021 Series) - Test 13 (31 May 2021)-img/p2_phy_s5_q10_1.png}{101.3}{0}
\qfigdef{p2_chem_s1_q03_1}{JEE Advanced (2021 Series) - Test 13 (31 May 2021)-img/p2_chem_s1_q03_1.png}{175.9}{0}
\qfigdef{p2_chem_s2_q03_A}{JEE Advanced (2021 Series) - Test 13 (31 May 2021)-img/p2_chem_s2_q03_A.png}{46.5}{0}
\qfigdef{p2_chem_s2_q03_2}{JEE Advanced (2021 Series) - Test 13 (31 May 2021)-img/p2_chem_s2_q03_2.png}{35.5}{0}
\qfigdef{p2_chem_s2_q03_D}{JEE Advanced (2021 Series) - Test 13 (31 May 2021)-img/p2_chem_s2_q03_D.png}{47.0}{0}
\qfigdef{p2_chem_s4_q03_1}{JEE Advanced (2021 Series) - Test 13 (31 May 2021)-img/p2_chem_s4_q03_1.png}{62.5}{0}
\qfigdef{p2_chem_s5_q02_1}{JEE Advanced (2021 Series) - Test 13 (31 May 2021)-img/p2_chem_s5_q02_1.png}{237.8}{0}
\qfigdef{p2_chem_s5_q03_1}{JEE Advanced (2021 Series) - Test 13 (31 May 2021)-img/p2_chem_s5_q03_1.png}{282.4}{1}
\qfigdef{p2_chem_s5_q09_1}{JEE Advanced (2021 Series) - Test 13 (31 May 2021)-img/p2_chem_s5_q09_1.png}{300.7}{1}
\qfigdef{p2_chem_s5_q10_1}{JEE Advanced (2021 Series) - Test 13 (31 May 2021)-img/p2_chem_s5_q10_1.png}{82.0}{0}

\begin{document}
\begin{center}{\Large\bfseries JEE Advanced (2021 Series) --- Test 13 (31 May 2021)}\end{center}
\begin{multicols}{2}
\papersec{1}
\partsec{1}{PHYSICS}
\sect{SEC--1}
\qn{1} A room air conditioner is modelled as a heat engine run in reverse. An amount of heat $Q_{in}$ is absorbed from the room at a temperature $T_L$ by the cooling coils which contains a working gas; this gas is compressed adiabatically to a temperature $T_H$; then the gas is compressed isothermally in a coil outside of the house, giving off an amount of heat $Q_{out}$; the gas expands adiabatically back to a temperature $T_L$ and finally expanded isothermally to original state and the cycle is repeated. An amount of energy $W$ is input into the system in every cycle through an electric pump. This model describes the air conditioner with the best possible efficiency.
\qfig{p1_phy_s1_q01_1}
The outside air temperature is $T_H$ and the inside air temperature is $T_L$. The air conditioner unit consumes electric power $P$. Assume that the air is dry sufficiently so that no condensation of water occurs in the cooling coils of the air conditioner. Which of the following may be the rate at which heat is removed from the room?
\begin{qoptions}
\opt{A}{$\dfrac P2\left(\dfrac{T_L}{T_H-T_L}\right)$}
\opt{B}{$\dfrac{3P}4\left(\dfrac{T_L}{T_H-T_L}\right)$}
\opt{C}{$P\left(\dfrac{T_L}{T_H-T_L}\right)$}
\opt{D}{$\dfrac{3P}2\left(\dfrac{T_L}{T_H-T_L}\right)$}
\end{qoptions}
\qn{2} A conducting medium is shaped in a form of a quarter of an annulus of radii $b$ and $a$ $(b>a)$ and thickness $t$. Specific resistance of the medium is $\rho$. The correct option(s) is/are
\qfig{p1_phy_s1_q02_1}
\begin{qoptions}
\opt{A}{Resistance between faces $LMNO$ and $PQRS$ is $\dfrac{2\rho}{\pi t}\ell n\dfrac ba$.}
\opt{B}{Resistance between faces $LMRS$ and $ONQP$ is $\dfrac{\pi\rho}{2t\ell n\left(\frac ba\right)}$.}
\opt{C}{Resistance between faces $MRQN$ and $LSPO$ is $\dfrac{4t\rho}{b^2-a^2}$.}
\opt{D}{Between faces $LMNO$ \& $PQRS$ and $LMRS$ \& $ONQP$, resistance is same.}
\end{qoptions}
\qn{3} Two resistors, one inductor and a capacitor are connected in a circuit with a battery and switch as shown in figure. The charge on capacitor and current in inductor is zero at $t=0$. Now switch is closed at $t=0$, then choose the correct option(s).
\qfig{p1_phy_s1_q03_1}
\begin{qoptions}
\opt{A}{The ratio of maximum current in the branch of inductor to the branch of capacitor is 200.}
\opt{B}{At $t=\dfrac{\ell n2}8$, the current in the branch of capacitor is 200 times the current in the branch of inductor.}
\opt{C}{At $t=\dfrac18\ell n\left(\dfrac54\right)$, the potential drop across inductor is 4 times the potential drop across capacitor.}
\opt{D}{The potential drop across inductor is maximum at $t=0$.}
\end{qoptions}
\qn{4} A small disc of mass $m$ is released on a parabolic curve in a vertical plane such that gravity acts along negative $y$-axis. The equation of parabolic curve is $x^2=\dfrac{2a}{\sqrt3}y$, where `$a$' is a positive constant. Frictional force between disc and curve are sufficient for pure rolling. When disc is reached at $x=a$ then choose the correct option(s).
\qfig{p1_phy_s1_q04_1}
\begin{qoptions}
\opt{A}{acceleration of disc along the trajectory is $\sqrt3g$.}
\opt{B}{acceleration of disc along the trajectory is $\dfrac g{\sqrt3}$.}
\opt{C}{frictional force between disc and curve is $\dfrac{mg}{2\sqrt3}$.}
\opt{D}{frictional force between disc and curve is $\dfrac{mg}{\sqrt3}$.}
\end{qoptions}
\qn{5} An object $O$ is kept infront of a converging lens of focal length 30cm behind which there is a plane mirror at 15cm from the lens.
\qfig{p1_phy_s1_q05_1}
\begin{qoptions}
\opt{A}{the final image is formed at 60cm from the lens towards right of it}
\opt{B}{the final image is at 60cm from lens towards left of it.}
\opt{C}{the final image is real.}
\opt{D}{the final image is virtual.}
\end{qoptions}
\qn{6} A light string is wrapped on a cylindrical shell and a fraction of length of string is unwrapped. A particle of mass $2M$ is attached on another end of string as shown. The system is kept in a vertical plane and cylinder can freely rotate about the axis of cylinder. Particle is released as shown in figure. Assuming there is no slipping between cylinder and string. Choose the correct option(s).
\qfig{p1_phy_s1_q06_1}
\begin{qoptions}
\opt{A}{The angular velocity of cylinder just after string taut is $\dfrac83\sqrt{\dfrac gR}$.}
\opt{B}{The angular velocity of cylinder just after string taut is $\dfrac{16}3\sqrt{\dfrac gR}$.}
\opt{C}{Velocity of particle just after string taut is $\dfrac83\sqrt{gR}$.}
\opt{D}{Velocity of particle just after string taut is $\dfrac{16}3\sqrt{gR}$.}
\end{qoptions}
\qn{7} A rod of length $2L$ and cross section area $A_0$ is kept in a liquid of density $\rho$ and pivoted at point $O$ in a vertical plane. The rod makes angle $30^\circ$ with the vertical. The mass density of rod is linearly increases from $O$ on both side as given by equation $\lambda=\lambda_0\left(1+\dfrac xL\right)$, where $\lambda_0$ is constant and $x$ is distance from point $O$ on either side. Rod is held in vertical plane by applying a force at end $A$ in horizontal direction. Then
\qfig{p1_phy_s1_q07_1}
\begin{qoptions}
\opt{A}{Centre of mass of rod $OA$ is at a distance $\dfrac{4L}9$ from point $A$.}
\opt{B}{Mass of rod $AB$ is $\dfrac{3\lambda_0L}2$.}
\opt{C}{Applied force at end $A$ is $\dfrac{17}{162\sqrt3}LA_0\rho g$.}
\opt{D}{Buoyant force will act at the centre of mass of rod inside the liquid.}
\end{qoptions}
\sect{SEC--2}
\paraq{1}{3}
In the Column-1, we have certain modification of YDSE apparatus/original YDSE apparatus. Column-2 gives the values of path difference and Column-3 gives resultant intensity at point $P$. Given : $d=1$mm, $D=1$m, $\lambda=500$ nm for all cases.
\end{multicols}
\begin{center}\small\setlength{\tabcolsep}{4pt}
\begin{tabular}{|p{0.9cm}|p{5.6cm}|p{0.9cm}|p{2.6cm}|p{0.9cm}|p{1.6cm}|}\hline
 & \textbf{Column-1} & & \textbf{Column-2} Path difference at $P$ & & \textbf{Column-3}\\\hline
(I) & \qfigi{p1_phy_s2_pre_1}\rule{0pt}{50pt}\newline Intensity due to either slit on screen $=I_0$. $P$ is centre of screen & (i) & Zero & (P) & $I_0$\\\hline
(II) & \qfigi{p1_phy_s2_pre_2}\rule{0pt}{44pt}\newline Slit 1 is of width $w$ and slit 2 is of width $2w$. Intensity due to slit 1 on screen $=I_0$. $P$ is centre of screen. & (ii) & 0.02 mm & (Q) & $4I_0$\\\hline
(III) & \qfigi{p1_phy_s2_pre_3}\rule{0pt}{44pt}\newline both slits are of same width. Slab of $\mu=1.5$ and thickness $0.5\ \mu$m placed in front of upper slit. $P$ is centre of screen. Intensity due to either slit on screen is $I_0$. & (iii) & $\dfrac16\ \mu$m & (R) & 0\\\hline
(IV) & \qfigi{p1_phy_s2_pre_4}\rule{0pt}{50pt}\newline $P$ is centre of screen. Intensity due to either slit on screen is $I_0$. Two slabs of equal width are placed in front of both the slits as shown. & (iv) & $0.25\ \mu$m & (S) & $5.8\,I_0$\\\hline
\end{tabular}
\end{center}
\begin{multicols}{2}
\qn{1} Which of the following is correct regarding minimum path difference?
\begin{qoptions}
\opt{A}{(II) (i) (S)}
\opt{B}{(I) (ii) (P)}
\opt{C}{(II) (iii) (Q)}
\opt{D}{(III) (i) (R)}
\end{qoptions}
\qn{2} Which of the following is correct regarding maximum intensity ?
\begin{qoptions}
\opt{A}{(IV) (iii) (Q)}
\opt{B}{(I) (ii) (Q)}
\opt{C}{(I) (ii) (P)}
\opt{D}{(III) (i) (R)}
\end{qoptions}
\qn{3} Which of the following is correct regarding minimum intensity ?
\begin{qoptions}
\opt{A}{(II) (i) (S)}
\opt{B}{(III) (iv) (R)}
\opt{C}{(I) (ii) (R)}
\opt{D}{(II) (iii) (R)}
\end{qoptions}
\paraq{4}{6}
A sphere, a cylinder, a spherical shell and a cylindrical shell are kept on a horizontal smooth surface. Two forces $F_1$ and $F_2$ are acting on each object. $F_1$ passes through centre of mass and $F_2$ acts on the top most point. All objects starts from rest on a smooth horizontal surface.
\end{multicols}
\begin{center}\small\setlength{\tabcolsep}{4pt}
\begin{tabular}{|p{0.9cm}|p{5.4cm}|p{0.9cm}|p{2.6cm}|p{0.9cm}|p{3.2cm}|}\hline
 & \textbf{Column-I} & & \textbf{Column-II} & & \textbf{Column-III}\\\hline
(I) & Uniform solid sphere of mass $m$, radius $R$\newline \qfigi{p1_phy_s2_q03_1}\rule{0pt}{62pt} & (i) & $F_1=F_2=F$ & (P) & Object will roll without sliding\\\hline
(II) & Uniform solid cylinder of mass $m$, radius $R$\newline \qfigi{p1_phy_s2_q03_2}\rule{0pt}{60pt} & (ii) & $F_2=2F_1=2F$ & (Q) & Object will roll on the horizontal surface with $a_{cm}<R\alpha$\\\hline
(III) & Uniform thin spherical shell of mass $m$, radius $R$\newline \qfigi{p1_phy_s2_q03_3}\rule{0pt}{62pt} & (iii) & $F_1=0,\ F_2=F$ & (R) & Object will roll on the horizontal surface with $2R\alpha>a_{cm}>R\alpha$\\\hline
(IV) & Uniform thin cylinder shell of mass $m$, radius $R$\newline \qfigi{p1_phy_s2_q03_4}\rule{0pt}{62pt} & (iv) & $F_1=\dfrac{3F_2}2=\dfrac{3F}2$ & (S) & Object will roll on the horizontal surface with $a_{cm}\ge2R\alpha$\\\hline
\end{tabular}
\end{center}
\begin{multicols}{2}
\qn{4} Mark the correct statement.
\begin{qoptions}
\opt{A}{(I) (i) (P)}
\opt{B}{(II) (ii) (Q)}
\opt{C}{(III) (ii) (R)}
\opt{D}{(IV) (iii) (S)}
\end{qoptions}
\qn{5} Mark the incorrect statement.
\begin{qoptions}
\opt{A}{(I) (ii) (Q)}
\opt{B}{(II) (i) (P)}
\opt{C}{(III) (iv) (R)}
\opt{D}{(IV) (iii) (S)}
\end{qoptions}
\qn{6} From the given combination the kinetic energy for time $t$ will be maximum.
\begin{qoptions}
\opt{A}{(I) (ii) (Q)}
\opt{B}{(II) (iv) (R)}
\opt{C}{(III) (iii) (P)}
\opt{D}{(IV) (i) (S)}
\end{qoptions}
\sect{SEC--3}
\qn{1} A particle having specific charge $\sigma$ is projected in $xy$ plane with a speed $v$. There exists a uniform magnetic field in $z$-direction having a fixed magnitude $B_0$. The field is made to reverse its direction after every interval of $\dfrac{2\pi}{\sigma B_0}$. Calculate the maximum separation (in m) between two positions of the particle during its course of motion. (given $\dfrac v{\sigma B_0}=2$ metre) (neglect any other force including gravity throughout the motion)
\qn{2} A shell of radius 1m is coated with a thin layer of $\beta^-$ active material. It's initial charge is zero and initial number of active atoms is $\dfrac43\times10^{12}$. If half life of decay is 1 hr and all the electron are emitted with an energy of 1.44 keV, find the time (in hr.) after which charge on the sphere becomes constant. Neglect the time taken by electron to return back.
\qn{3} Half portion of a spherical capacitor is filled with a dielectric of dielectric constant $k=2$ and conductivity $\sigma$. The charge given to spherical capacitor is $Q_0$. Due to the conductivity of dielectric charge leaks. The time constant for the discharge circuit is $\dfrac{n\epsilon_0}{\sigma}$. Find the value of $n$.
\qfig{p1_phy_s3_q03_1}
\qn{4} Figure shows a soap film formed between two square figures made of a uniform wire. The bigger square is held while keeping it in a horizontal plane and the smaller square is slowly allowed to drop vertically. It reaches an equilibrium state after dropping a height $h$. Let surface tension of soap $=T$. Mass per unit length of wire $=\lambda$. Acceleration due to gravity $=g$. Given that $h=\dfrac{n\lambda ga}{2\sqrt{4T^2-\lambda^2g^2}}$; find the integer value $n$.
\qfig{p1_phy_s3_q04_1}
\qn{5} A pipe 20 cm long is closed at one end. Which harmonic mode of the pipe will be resonantly exited by a source of 1237.5 Hz. Sound velocity in air 330 m/s.
\partsec{2}{CHEMISTRY}
\sect{SEC--1}
\qn{1} Choose the correct statements :
\begin{qoptions}
\opt{A}{Temperature, enthalpy and entropy are state functions}
\opt{B}{For reversible and irreversible both isothermal expansion of an ideal gas, change in internal energy is zero}
\opt{C}{For a reaction in which $\Delta n_g=0$, entropy change is not always zero}
\opt{D}{The entropy change of an ideal gas due to its reversible isothermal expansion is equal to $2.303\,R\log_{10}\dfrac{P_1}{P_2}$ [where $P_1$ \& $P_2$ are initial and final pressures of gas]}
\end{qoptions}
\qn{2} Carefully study the following 2 cases and choose the correct statements followed ?\\
Case \# I. All H atoms in ethene are substituted differently without significantly distorting the original geoemetry.\\
Case \# II. All H atoms in ethane are substituted differently without significantly distorting the original geoemetry.\\
Choose the correct statements(s) from the following :
\begin{qoptions}
\opt{A}{Diastereomers will be produced only in Case I , not in Case II.}
\opt{B}{Enantiomers will be produced only in Case II, not in Case I}
\opt{C}{Mixing equal weights of one dextro and one laevo molecule under Case II will necessarily yield an optically inactive mixture}
\opt{D}{Total number of space isomers possible under both cases $=86$}
\end{qoptions}
\qfullbegin
\qn{3} Consider the following reaction
\qfig{p1_chem_s1_q03_1}
\qfullsplit
Which of the following is/are correct regarding the above reaction.
\begin{qoptions}
\opt{A}{I and II have same value of EAN}
\opt{B}{CO formed in the reaction can be used to increase rate of reaction in Haber's process.}
\opt{C}{I and II both consist of trihapto ligands.}
\opt{D}{The above reaction is redox reaction}
\end{qoptions}
\qfullend
\qn{4} Inorganic benzene surpasses organic benzene
\begin{qoptions}
\opt{A}{Melting point}
\opt{B}{B--H bond length data as compare to C--H bond length data}
\opt{C}{Molar mass}
\opt{D}{Total number of electrons in the molecule}
\end{qoptions}
\qn{5} If equatorial plane in PCl$_5$ molecule is the $X$--$Y$ plane, the orbitals hybridizing to produce axial bonds will be -
\begin{qoptions}
\opt{A}{$p_z$}
\opt{B}{$d_{yz}$}
\opt{C}{$p_y$}
\opt{D}{$d_{z^2}$}
\end{qoptions}
\qn{6} Which of the following changes can be experimentally realized in the laboratory:
\begin{qoptions}
\opt{A}{Combustion of Mg ribbon in CO$_2$}
\opt{B}{Reduction of FeCl$_3$ (aq.) on adding Zn granules}
\opt{C}{Direct heating of hydrated MgCl$_2$ to get anhydrous MgCl$_2$}
\opt{D}{Evolution of brown gas on heating LiNO$_3$}
\end{qoptions}
\qn{7} An optically active aldotetrose which retains its optical activity with aqueous Br$_2$ but loses it with conc.HNO$_3$ has the following structure:
\begin{qoptions}
\opt{A}{\qfigi{p1_chem_s1_q07_A}}
\opt{B}{\qfigi{p1_chem_s1_q07_B}}
\opt{C}{\qfigi{p1_chem_s1_q07_C}}
\opt{D}{\qfigi{p1_chem_s1_q07_D}}
\end{qoptions}
\sect{SEC--2}
\end{multicols}
\begin{center}\small\setlength{\tabcolsep}{4pt}
\begin{tabular}{|p{0.8cm}|p{3cm}|p{0.9cm}|p{5.2cm}|p{0.9cm}|p{4.6cm}|}\hline
 & \textbf{Column-I} (Polymer) & & \textbf{Column-II} (Monomer) & & \textbf{Column-III} (Properties of monomer)\\\hline
(I) & Polystyrene & (i) & Caprolactum & (P) & Resonance stabilization\\\hline
(II) & Poly Vinyl Chloride & (ii) & PhCH $=$ CH$_2$ & (Q) & 1,4 addition of HBr (mainly) at 40$^\circ$C\\\hline
(III) & Natural rubber & (iii) & CH$_2$ $=$ CHCl & (R) & 1,2 addition of HBr at room temp. (mainly and as per traditional Markonikoff's rule)\\\hline
(IV) & Nylon-6 & (iv) & CH$_2$ $=$ C(CH$_3$) $-$ CH $=$ CH$_2$ & (S) & CO$_2$(g) evolution with KMnO$_4$/H$^+$\\\hline
\end{tabular}
\end{center}
\begin{multicols}{2}
\qn{1} Which of the following combination is correct?
\begin{qoptions}
\opt{A}{(III) (ii) (R)}
\opt{B}{(II) (i) (P)}
\opt{C}{(IV) (i) (P)}
\opt{D}{(I) (iii) (S)}
\end{qoptions}
\qn{2} Which of the following combination is incorrect?
\begin{qoptions}
\opt{A}{(I) (ii) (P)}
\opt{B}{(I) (ii) (R)}
\opt{C}{(III) (iv) (P)}
\opt{D}{(II) (iv) (R)}
\end{qoptions}
\qn{3} Which of the following combination is correct?
\begin{qoptions}
\opt{A}{(II) (ii) (P)}
\opt{B}{(II) (iii) (R)}
\opt{C}{(III) (iii) (P)}
\opt{D}{(IV) (ii) (P)}
\end{qoptions}
\end{multicols}
\begin{center}\small\setlength{\tabcolsep}{4pt}
\begin{tabular}{|p{0.8cm}|p{3.8cm}|p{0.9cm}|p{4.6cm}|p{0.9cm}|p{3.6cm}|}\hline
 & \textbf{Column-I} & & \textbf{Column-II} & & \textbf{Column-III}\\\hline
(I) & Peptization & (i) & Colloid preparation technique & (P) & Sky in blue\\\hline
(II) & Delta formation & (ii) & Colloid destabilization & (Q) & Gold number\\\hline
(III) & Tyndall effect & (iii) & Colloid stabilization & (R) & Coagulation\\\hline
(IV) & Protection power of colloids & (iv) & Optical property of colloid & (S) & Adsorption of common ion\\\hline
\end{tabular}
\end{center}
\begin{multicols}{2}
\qn{4} Which of the following combination is correct?
\begin{qoptions}
\opt{A}{(I) (i) (S)}
\opt{B}{(II) (i) (R)}
\opt{C}{(III) (ii) (S)}
\opt{D}{(I) (ii) (S)}
\end{qoptions}
\qn{5} Which of the following combination is correct?
\begin{qoptions}
\opt{A}{(II) (i) (S)}
\opt{B}{(II) (ii) (R)}
\opt{C}{(III) (ii) (S)}
\opt{D}{(IV) (iii) (R)}
\end{qoptions}
\qn{6} Which of the following combination is INCORRECT?
\begin{qoptions}
\opt{A}{(III) (iv) (P)}
\opt{B}{(II) (ii) (R)}
\opt{C}{(IV) (iii) (Q)}
\opt{D}{(I) (ii) (Q)}
\end{qoptions}
\sect{SEC--3}
\qn{1} Ammonium carbamate is heated at $227^\circ$C in a closed vessel of volumn 8210 ml containing 0.2 mole of Argon gas. The manometer attached to the vessel shows a pressure of 4 atm. Find $K_p$ (in atm$^3$) for the decomposition of ammonium carbamate.
\[\mathrm{NH_2COONH_4\,(s)\rightleftharpoons 2NH_3(g)+CO_2(g)}\]
\qn{2} How many gases/vapours are involved either as reactants or products in Serpeck's process used for concentration of bauxite.
\qn{3}
\qfig{p1_chem_s3_q03_1}
When the above molecule is reduced with excess LiAlH$_4$ followed by dilute acid hydrolysis, how many organic products can be isolated by fractional distillation of product mixture.
\qn{4} The pH at which a $0.1$ M Al$^{+3}$ solution is $99.99\%$ precipitated is :\\
($K_{sp}$ Al(OH)$_3=1\times10^{-20}$ at $25^\circ$C.)
\qn{5} No. of reagents which on reaction with thiosulphate solution can yield a precipitate or turbidity.\\
KI$_3$, BaCl$_2$, CaCl$_2$, Pb(NO$_3$)$_2$, KCN, HCl, AgNO$_3$
\partsec{3}{MATHEMATICS}
\sect{SEC--1}
\qn{1} Let $P=\displaystyle\sum_{n=1}^{10}n\left\{\dfrac{1^2}{1+n}+\dfrac{2^2}{2+n}+\ldots+\dfrac{10^2}{10+n}\right\}$, then which of the following is/are correct ?
\begin{qoptions}
\opt{A}{$2P$ is perfect square}
\opt{B}{$2P$ is divisible by $11^2$}
\opt{C}{Number of divisors of $2P$ is 9}
\opt{D}{The remainder when $2P$ is divided by 6 is 1}
\end{qoptions}
\qn{2} Define a function $f(x)=\begin{cases}\sin x, & \text{if }x\text{ is irrational}\\ \tan x, & \text{if }x\text{ is rational}\end{cases}$. Then $f(x)$ is
\begin{qoptions}
\opt{A}{not continuous at any point in the interval $(-\pi,\pi)$}
\opt{B}{continuous at some point(s) , but not differentiable at any point in the interval $(-\pi,\pi)$}
\opt{C}{continuous at some point(s) in the interval $(-\pi,\pi)$, but not differentiable at some of these point(s)}
\opt{D}{continuous at some point(s) in the interval $(-\pi,\pi)$ and differentiable at all these point(s)}
\end{qoptions}
\qn{3} Let $A=\begin{bmatrix}1&2\\-a^2&a+2\end{bmatrix}$ and $B,C$ are $2\times2$ matrices such that $ABC=\begin{bmatrix}4&2\\1&3\end{bmatrix}$. Let $(BC)^T=[p_{ij}]_{2\times2}$ and $p_{21}=-1$, then the value of `$a$' can be-
\begin{qoptions}
\opt{A}{0}
\opt{B}{$-\dfrac32$}
\opt{C}{$\dfrac12$}
\opt{D}{$-\dfrac12$}
\end{qoptions}
\qn{4} A ray of light coming along the line $y=x+1$ from positive $x$-direction strikes a concave mirror whose intersection with $xy$-plane is the parabola $(y-x)^2=16(x+y)$. Then which of the following is/are correct statement(s)
\begin{qoptions}
\opt{A}{reflected ray passes through $(2,2)$}
\opt{B}{slope of reflected ray is $\dfrac{47}{79}$}
\opt{C}{reflected ray passes through $(1,1)$}
\opt{D}{slope of normal at point of contact is $\dfrac79$}
\end{qoptions}
\qn{5} Let $L_1:x+y=1,\ z+2y=2$ and $L_2:x-y=1,\ z-2y=2$ be two straight lines, then which of the following is/are correct statement(s)?
\begin{qoptions}
\opt{A}{the lines $L_1$ and $L_2$ are coplanar}
\opt{B}{the lines $L_1$ and $L_2$ are non-coplanar}
\opt{C}{on the line $\vec r=4\hat j+\lambda(\hat i+2\hat k)$, $\lambda\in R$, there exist infinite number of points such that an infinite number of lines passing through each of these points and intersecting both the lines $L_1$ and $L_2$ can be drawn}
\opt{D}{there exists a unique line passing through the point $A(1,\mu,3)$ which intersect both lines $L_1$ and $L_2$ $\forall\ \mu\in R$}
\end{qoptions}
\qn{6} Let $I=\displaystyle\int_1^2\dfrac{e^x+e^{4/x}}{x}dx$, then which of the following is/are correct statement(s) ?
\begin{qoptions}
\opt{A}{$2I=\displaystyle\int_1^4\dfrac{e^x+e^{4/x}}{x}dx$}
\opt{B}{$4I=\displaystyle\int_1^4\dfrac{e^x+e^{4/x}}{x}dx$}
\opt{C}{$I>(e^2-e)$}
\opt{D}{$I<(e^2-e)$}
\end{qoptions}
\qn{7} Let $S_n=\displaystyle\sum_{r=1}^n\cot^{-1}\left(\dfrac{9+r(r-1)}3\right)$, then which of the following is/are correct statement(s) ?
\begin{qoptions}
\opt{A}{$S_n=\sin^{-1}\left(\dfrac3{\sqrt{n^2-9}}\right)$}
\opt{B}{$S_n=\sin^{-1}\left(\dfrac n{\sqrt{n^2+9}}\right)$}
\opt{C}{$S_n>\dfrac n3$}
\opt{D}{$S_n<\dfrac n3$}
\end{qoptions}
\sect{SEC--2}
\paraq{1}{3}
Let $f(x)=\dfrac{(x-a)(x-b)}{(x-c)(x-d)}$ where $a<b<c<d$. Column-I: contains information about real roots of the equation; Column-II: contains information about limit of $f(x)$; Column-III: contains information about increasing interval and decreasing interval of $f(x)$ and $f'(x)$ and information about local minimum value and local maximum value of $f(x)$
\end{multicols}
\begin{center}\small\setlength{\tabcolsep}{4pt}
\begin{tabular}{|p{0.8cm}|p{4.4cm}|p{0.8cm}|p{3.4cm}|p{0.8cm}|p{4.2cm}|}\hline
 & \textbf{Column-I} & & \textbf{Column-II} & & \textbf{Column-III}\\\hline
(I) & $f(x)=1$ has exactly one real root & (i) & $\displaystyle\lim_{x\to c^-}f(x)=-\infty$ & (P) & $f(x)$ is increasing in the interval $(b,c)$\\\hline
(II) & $f(x)=2$ has exactly two real root & (ii) & $\displaystyle\lim_{x\to d^+}f(x)=\infty$ & (Q) & $f(x)$ is decreasing in the interval $(d,\infty)$\\\hline
(III) & $f(x)=-10$ has one real root & (iii) & $\displaystyle\lim_{x\to d^-}f(x)=-\infty$ & (R) & $f'(x)$ is increasing in the interval $(b,c)$\\\hline
(IV) & $f(x)=10$ has two real root & (iv) & $\displaystyle\lim_{x\to\infty}f(x)=1$ & (S) & local minimum value of $f(x)$ is smaller than local maximum value of $f(x)$\\\hline
\end{tabular}
\end{center}
\begin{multicols}{2}
\qn{1} Which of the following options is the only CORRECT combination ?
\begin{qoptions}
\opt{A}{(I) (ii) (S)}
\opt{B}{(II) (iii) (Q)}
\opt{C}{(I) (i) (R)}
\opt{D}{(III) (iv) (S)}
\end{qoptions}
\qn{2} Which of the following options is the only INCORRECT combination ?
\begin{qoptions}
\opt{A}{(II) (ii) (P)}
\opt{B}{(III) (i) (R)}
\opt{C}{(IV) (iv) (P)}
\opt{D}{(I) (iii) (Q)}
\end{qoptions}
\qn{3} Which of the following options is the only CORRECT combination ?
\begin{qoptions}
\opt{A}{(III) (iv) (P)}
\opt{B}{(I) (ii) (R)}
\opt{C}{(IV) (i) (S)}
\opt{D}{(II) (iv) (S)}
\end{qoptions}
\paraq{4}{6}
Let $f(x)$ be a double differentiable function with continuous second derivative defined on the interval $[0,1]$ satisfying the equation $\displaystyle\int_0^x\sqrt{1-(f''(t))^2}\,dt=\int_0^xf'(t)\,dt\ \forall x\in[0,1]$ and $f(0)=f'(0)=0$. Column-I: contains the information about interval of value of $f(x)$; Column-II: contains information about value of $f(x)$ ; Column-III: contains information about interval of value of definite integral of $f(x)$; Match the following Column(s)
\end{multicols}
\begin{center}\small\setlength{\tabcolsep}{4pt}
\begin{tabular}{|p{0.8cm}|p{3cm}|p{0.9cm}|p{4.2cm}|p{0.9cm}|p{3.2cm}|}\hline
 & \textbf{Column-I} & & \textbf{Column-II} & & \textbf{Column-III}\\\hline
(I) & $f\left(\dfrac12\right)<\dfrac18$ & (i) & $f\left(\dfrac13\right)=\dfrac1{18}$ & (P) & $\displaystyle\int_0^1f(x)\,dx>1$\\\hline
(II) & $f\left(\dfrac12\right)>\dfrac18$ & (ii) & $f\left(\dfrac14\right)=\dfrac18$ & (Q) & $\displaystyle\int_0^1f(x)\,dx>\dfrac16$\\\hline
(III) & $f\left(\dfrac13\right)<\dfrac1{18}$ & (iii) & $f\left(\dfrac16\right)=2\sin^2\left(\dfrac1{12}\right)$ & (R) & $\displaystyle\int_0^1f(x)\,dx<\dfrac16$\\\hline
(IV) & $f\left(\dfrac14\right)<\dfrac14$ & (iv) & $f\left(\dfrac18\right)=2\sin^2\left(\dfrac1{16}\right)$ & (S) & $\displaystyle\int_0^1f(x)\,dx<\dfrac14$\\\hline
\end{tabular}
\end{center}
\begin{multicols}{2}
\qn{4} Which of the following options is the only CORRECT combination ?
\begin{qoptions}
\opt{A}{(I) (ii) (Q)}
\opt{B}{(II) (i) (S)}
\opt{C}{(IV) (ii) (Q)}
\opt{D}{(III) (iii) (R)}
\end{qoptions}
\qn{5} Which of the following options is the only CORRECT combination ?
\begin{qoptions}
\opt{A}{(I) (iv) (Q)}
\opt{B}{(II) (iii) (S)}
\opt{C}{(II) (iv) (P)}
\opt{D}{(IV) (iii) (S)}
\end{qoptions}
\qn{6} Which of the following options is the only CORRECT combination ?
\begin{qoptions}
\opt{A}{(III) (ii) (Q)}
\opt{B}{(IV) (ii) (P)}
\opt{C}{(I) (iii) (R)}
\opt{D}{(III) (ii) (S)}
\end{qoptions}
\sect{SEC--3}
\qn{1} The number of $x$ in $(0,2\pi)$ satisfying the equation $1+\sin2x=\sin x+\sin^2x$ is
\qn{2} If $z_1,z_2$ are complex numbers, then the maximum value of $\dfrac{(z_1\bar z_2+\bar z_1z_2)+(z_1z_2+\bar z_1\bar z_2)}{|z_1z_2|}$ is equal to
\qn{3} Let $PQ$ be a focal chord of $\dfrac{x^2}9+\dfrac{y^2}4=1$ passing through one focus $S$ and $S'$ be its other focus, then the value of semi-perimeter of $\Delta S'PQ$ is
\qn{4} Let $f(x)=ax^2+bx+c$, where $a\ne0$, $a,b,c$ are integers and $f(1)=1$, $6<f(3)<8$ and $18<f(5)<22$. Then the number of solutions of equation $f(x)=e^x$ is
\qn{5} Let $L=\displaystyle\lim_{x\to\infty}\left(\left(\dfrac{x+1}{x-1}\right)^x-e^2\right)x^2$, then the value of $\dfrac{9L}{2e^2}$ is equal to
\papersec{2}
\partsec{1}{PHYSICS}
\sect{SEC--1}
\qn{1} A practically very small source ($S$) of light is kept in front of a plane mirror. But due to imperfect manufacturing, there is a small spherical bulge (radius of curvature 80 mm). The situation is as shown in the figure. You can assume that principal axis of bulge is perpendicular to the plane mirror. Find the separation between the images formed.
\qfig{p2_phy_s1_q01_1}
\begin{qoptions}
\opt{A}{10.025 mm}
\opt{B}{20.025 mm}
\opt{C}{40.025 mm}
\opt{D}{80.025 mm}
\end{qoptions}
\qn{2} A light ray incident along vector $2\hat i+4\hat j+\sqrt5\,\hat k$ strikes on the $x$-$z$ plane from medium of refractive index $\sqrt3$ and enters into medium II of refractive index is $\mu_2$. The value of $\mu_2$ for which the value of angle of refraction becomes $90^\circ$, is
\begin{qoptions}
\opt{A}{$\dfrac{4\sqrt3}{5}$}
\opt{B}{$\dfrac{3\sqrt3}{5}$}
\opt{C}{$\dfrac{2\sqrt3}{5}$}
\opt{D}{$\dfrac{\sqrt3}{5}$}
\end{qoptions}
\qn{3} A solid sphere has density $\rho$ in the top half and $2\rho$ in the bottom half. It is kept in a uniform fluid of density $3\rho$ with a hinge at the bottom (as shown) such that it can freely perform oscillations about any horizontal axis through the hinge (Radius of sphere is $R$). Find the period of small oscillations if all the points of the sphere are to oscillate in their respective planes parallel to the plane $x=-\sqrt3z$ (shown).
\qfig{p2_phy_s1_q03_B}
\begin{qoptions}
\opt{A}{$2\pi\sqrt{\dfrac{64R}{35g}}$}
\opt{B}{$2\pi\sqrt{\dfrac{46R}{45g}}$}
\opt{C}{$2\pi\sqrt{\dfrac{32R}{35g}}$}
\opt{D}{$2\pi\sqrt{\dfrac{24R}{35g}}$}
\end{qoptions}
\sect{SEC--2}
\qn{1} Two balls $A$ \& $B$ of mass $m_1$ and $m_2$ are kept on a horizontal smooth surface. $A$ is given a velocity towards $B$ so that they perform head on collision
\begin{qoptions}
\opt{A}{If $m_1=m_2$ and collision is elastic $A$ stops and $B$ moves with the velocity of $A$ after collision}
\opt{B}{Impulses on $A$ and $B$ during collision are same in magnitude whether it is elastic or inelastic collision}
\opt{C}{Impulses on $A$ and $B$ during collision are equal in magnitude only if $m_1=m_2$}
\opt{D}{If $m_1>m_2$, though impulses are same in magnitude velocity of $A$ is still in the same direction.}
\end{qoptions}
\qn{2} A uniform rod of length $\ell$ and mass $m$ is rotating about a fixed axis perpendicular to the paper ($xy$-plane) and passing through its one of the end $A$. At the instant shown in the figure rod's angular velocity is $\vec\omega$ and velocity of the centre of mass of the rod is $\vec v_{cm}$. The moment of inertia of the rod about an axis passing through its centre of mass and parallel to $\hat k$ is $I_{CM}$ and the moment of inertia of the rod about an axis passing through point $A$ and parallel to $\hat k$ is $I_A$. $\vec L_A$ and $K$ represent the angular momentum of the rod about point $A$ and kinetic energy of rod respectively.
\qfig{p2_phy_s2_q02_1}
\begin{qoptions}
\opt{A}{$\vec L_A=I_A\vec\omega=\left(I_{cm}+\dfrac{m\ell^2}{4}\right)\vec\omega,\ K=\dfrac12mv_{cm}^2+\dfrac12I_{cm}\omega^2$}
\opt{B}{$\vec L_A=I_A\vec\omega=m\vec r_{cm}\times\vec v_{cm}\,\vec\omega,\ K=\dfrac12mv_{cm}^2+\dfrac12I_{cm}\omega^2$}
\opt{C}{$\vec L_A=I_A\vec\omega=\left(I_{cm}+\dfrac{m\ell^2}{4}\right)\vec\omega,\ K=\dfrac12I_A\omega^2$}
\opt{D}{$\vec L_A=I_A\vec\omega=m\vec r_{cm}\times\vec v_{cm}+I_{cm}\vec\omega,\ K=\dfrac12I_A\omega^2$}
\end{qoptions}
\qn{3} In young's experiment, the upper slit is covered by a thin glass plate of refractive index $\dfrac43$ and thickness $9\lambda$, where $\lambda$ is the wavelength of light used in the experiment. The lower slit is also covered by another glass plate of thickness $2\lambda$ \& refractive index $\dfrac32$, as shown in figure. If $I_0$ is the intensity at point $P$ due to slits $S_1$ \& $S_2$ each, then:
\qfig{p2_phy_s2_q03_1}
\begin{qoptions}
\opt{A}{Intensity at point $P$ is $4I_0$.}
\opt{B}{Two fringes have been shifted in upward direction after insertion of both the glass plates together.}
\opt{C}{Optical path difference between the waves from $S_1$ \& $S_2$ at point $P$ is $2\lambda$.}
\opt{D}{If the source $S$ is shifted upwards by a small distance $d_2$ then the fringe originally at $P$ after inserting the plates, shifts downwards by $D\left(\dfrac{d_2}{d_1}\right)$.}
\end{qoptions}
\qn{4} Choose the correct options regarding to the following A.C. circuit:
\qfig{p2_phy_s2_q04_1}
\begin{qoptions}
\opt{A}{Resonance occurs when $\omega=\dfrac1{\sqrt{LC}}$.}
\opt{B}{Resonance occurs when power factor is maximum.}
\opt{C}{Resonance occurs when reactive power of circuit becomes equal to zero.}
\opt{D}{Resonance occurs when $\omega=\sqrt{\dfrac1{LC}-\left(\dfrac RL\right)^2}$.}
\end{qoptions}
\qn{5} A block of mass $m=1$ kg is placed on a plank of mass $M=2$ kg, which is placed on a horizontal frictionless surface. There is no friction between the block and the plank. Block and plank are connected by a spring of spring constant 6 N/m as shown in the figure. An impulse $J=5$ newton-sec is applied on the plank. Possible magnitude of acceleration of the block is/are :
\qfig{p2_phy_s2_q05_1}
\begin{qoptions}
\opt{A}{1 m/s$^2$}
\opt{B}{3 m/s$^2$}
\opt{C}{5 m/s$^2$}
\opt{D}{7 m/s$^2$}
\end{qoptions}
\sect{SEC--3}
\paraq{1}{2}
Let us consider a diatomic gas whose molecules have the shape of a dumbbell. In this model, the center of mass of the molecule can translate in the $x$, $y$, and $z$ directions. In addition, the molecule can rotate about three mutually perpendicular axes. We can neglect the rotation about the $y$-axis because the moment of inertia $I_y$ and the rotational energy $\frac12I_y\omega^2$ about this axis are negligible compared with those associated with the $x$ and $z$ axes. (If the two atoms are taken to be point masses, then $I_y$ is identically zero). Thus, there are five degrees of freedom: three associated with the translational motion and two associated with the rotational motion. Because each degree of freedom contributes, on the average, $\frac12k_BT$ of energy per molecule, the total internal energy for a system of $N$ molecules is:
\begin{gather*}E_{int}=3N\left(\tfrac12k_BT\right)+2N\left(\tfrac12k_BT\right)\\=\tfrac52Nk_BT=\tfrac52nRT\ \ldots\text{(i)}\end{gather*}
We can use this result to find the molar specific heat at constant volume:
\qfig{p2_phy_s3_para1_2_1}
\[C_V=\frac1n\frac{dE_{int}}{dT}=\frac1n\frac d{dT}\left(\tfrac52nRT\right)=\tfrac52R\ \ldots\text{(ii)}\]
From equation (i) and (ii), we find that
\[C_P=C_V+R=\tfrac72R,\qquad \gamma=\frac{C_P}{C_V}=\frac{7/2R}{5/2R}=\frac75=1.40\]
In the vibratory model, the two atoms are joined by an imaginary spring. The vibrational motion adds two degrees of freedom, which correspond to the kinetic energy and the potential energy associated with vibrations along the length of the molecule about its centre of mass.
\qn{1} The root mean square angular velocity of a diatomic molecule (with each atom of mass $m$ and interatomic distance $a$) is given by:-
\begin{qoptions}
\opt{A}{$\sqrt{\dfrac{4k_BT}{ma^2}}$}
\opt{B}{$\sqrt{\dfrac{2k_BT}{ma^2}}$}
\opt{C}{$\sqrt{\dfrac{k_BT}{ma^2}}$}
\opt{D}{$\sqrt{\dfrac{k_BT}{2ma^2}}$}
\end{qoptions}
\qn{2} A diatomic molecule is moving without rotation or vibration with velocity $v_{rms}$ such that it is oriented along $x$-axis. It strikes a wall in $yz$-plane while moving in $+$ve $x$ direction. The spring constant can be assumed to be $K$ and time of collision is negligible. After all (elastic) collision are over:-
\begin{qoptions}
\opt{A}{The molecule is moving along $-$ve $x$-direction and oscillating about its centre of mass}
\opt{B}{The molecule is moving along $-$ve $x$-direction, but not oscillating about its centre of mass}
\opt{C}{The molecule is oscillating about it centre of mass but not moving at all.}
\opt{D}{The molecule is neither rotating nor moving at all.}
\end{qoptions}
\sect{SEC--4}
\qn{1} A lamp post of length $\ell$ and cross sectional area $S=10$ cm$^2$ is made of material with linear expansion coefficient $\alpha=1.1\times10^{-5}$ K$^{-1}$ and Young's modulus $Y=2.0\times10^{11}$ N/m$^2$. If temperature of lamp post is raised by $\Delta T=10$K, what mass (in kg) should be put on the lamp post, so that there is no change in the length of lamp post ($g=10$ m/s$^2$)
\qn{2} Two ideal solenoids of same dimensions. One is air cored with 600 turns while other is Aluminium cored with 200 turns (relative permeability of Aluminium is 3), are connected in a circuit as shown in the figure. The switch $S$ is closed at $t=0$. Find the ratio of potential difference across air-cored solenoid to that of Aluminium cored solenoid at any time $t$.
\qfig{p2_phy_s4_q02_1}
\qn{3} A cubical block of wood (density = 800 kg/m$^3$) and side length $\ell=50$ cm is floating in a rectangular tank partially filled with water (density = 1000 kg/m$^3$). Work done to press the block slowly so that it is just immersed in water is $W$ Joules (Area of the bottom of tank is 10000 cm$^2$) (take $g=10$ m/s$^2$). Find the value of $W$.
\sect{SEC--5}
\qn{1} It is known that the electrically electrified bodies can attract the uncharged body. For example, the charged comb attracts the uncharged pieces of paper. One can estimate the force of interaction between charged body and uncharged body. For the estimation of force, consider an electrified body, localised at point $A$, having charge $q$ and an uncharged body in form of conducting cylinder whose dimensions are considerably lower than distance from charge $q$ ($a\gg\ell$). The axis of cylinder is passing through point $A$. The volume of cylinder is $V$. If the force of interaction between them is $\dfrac{q^2}{\alpha\pi^2\varepsilon_0a^5}V$. Find the value of $\alpha$. (assume that no charge appears on the curved surface of cylinder)
\qfig{p2_phy_s5_q01_1}
\qn{2} The small particle of mass $m$ is given an initial high velocity in the horizontal plane and winds its cord around the fixed vertical shaft of radius 1 m. All motion occurs essentially in the horizontal plane. If the angular velocity of the cord is 0.8 rad/s when the distance from the particle to the tangency point is 5m., determine the angular velocity $\omega$(in rad/s) of the cord after it has turned through an angle 1 rad.
\qfig{p2_phy_s5_q02_1}
\qn{3} An inclined plane is located at angle $\alpha=53^\circ$ to the horizontal. There is a hole at point $B$ in inclined plane as shown in the figure. A particle is projected along the plane with speed $v_0$ at an angle $\beta=37^\circ$ to the horizontal in such a way that it get into the hole. Neglect any type of friction. Find the speed $v_0$ (in m/s) if $h=1$ m, $\ell=8$m and $\sqrt5\cong2.25$ .
\qfig{p2_phy_s5_q03_1}
\qn{4} A non conducting solid cylinder of infinite length having charge density $\rho$ and radius of cylinder is $5R$. If the flux passing through the surface $ABCD$ as shown in figure is $\dfrac{2\theta\rho R^3}{\varepsilon_0}$. Find the value of $\theta$.
\qfig{p2_phy_s5_q04_1}
\qn{5} The circuit shows a resistance $R=0.01\ \Omega$ and inductance $L=3$ mH connected to a conducting rod $PQ$ of length $\ell=2$m which can side on a perfectly conducting circular arc of radius $\ell$ with its centre at $P$. Assume that friction \& gravity are absent and a constant uniform magnetic field $B=0.1$ T exists as shown in the figure. At $t=0$, the circuit is switched on and simultaneously an external torque is applied on the rod so that it rotates about $P$ with a constant angular velocity $\omega=2$ rad/sec. Find the magnitude of this torque (in N-m) at $t=(0.3\ \ell n\,2)$ second
\qfig{p2_phy_s5_q05_1}
\qn{6} A meter stick connected with disc of same mass which is joined at one end of stick and disc can freely rotates about its joint. System swings in a vertical plane about a fixed horizontal axis passing through its one end, undergoes small oscillation of frequency `$f_0$'. If disc is removed from the stick, then its new frequency of small oscillation is $\sqrt{\dfrac4k}\,f_0$. Find the value of $k$
\qfig{p2_phy_s5_q06_1}
\qn{7} Two capacitors $C_1=C_2=\dfrac1{\pi^2}\times10^{-2}$F and inductor $L=2\times10^{-2}$ H are connected in series as shown in the figure. Initially charge on each capacitors are $4\sqrt3\,\mu$C. At $t=0$ switch $S_1$ is closed and $t=1/400$ sec, switch $S_2$ is also closed. The maximum charge on capacitor $C_2$ during $LC$ oscillation is $k\sqrt2\,\mu$C. Find $k$.
\qfig{p2_phy_s5_q07_1}
\qn{8} Two identical equi-concave lenses made of glass of reflected index 1.5 placed in contact has power $P$. When a liquid of refractive index $\mu$ is filled in the gap between the concave lenses, the power becomes of initial value $(P/3)$. The value of $\mu$ is $\dfrac K3$. Find the value of $K$.
\qn{9} A hopper drops gravel at a rate of 15 kg/s onto a conveyor belt moving at a constant speed 4m/s. If the output power must motor have in order to drive the conveyor belt is 48 $k$. Find the value of $k$.
\qn{10} A parallel plate capacitor (square plates) has a capacitance of 1 mF. It is completely filled with a solid dielectric (constant = 5) which is performing very small oscillations of amplitude 0.1\% of plate dimension with a period of 1 ms as shown: At all times, a battery of 10 V remains connected to the capacitor. If the amount of charge that crosses the battery in one cycle is $40\ \alpha$. Find the value of $\alpha$.
\qfig{p2_phy_s5_q10_1}
\partsec{2}{CHEMISTRY}
\sect{SEC--1}
\qn{1} Correct trend of first ionisation energy in group-13 is
\begin{qoptions}
\opt{A}{B $>$ Al $>$ Ga $>$ In $>$ Tl}
\opt{B}{B $>$ Al $>$ Ga $>$ Tl $>$ In}
\opt{C}{B $>$ Tl $>$ Ga $>$ Al $>$ In}
\opt{D}{B $>$ Ga $>$ Al $>$ In $>$ Tl}
\end{qoptions}
\qn{2} The self ionization constant for pure formic acid $K=[\mathrm{HCOOH_2^+}][\mathrm{HCOO^-}]$ has been estimated as $10^{-6}$ at room temperature. The density of formic acid is 1.15 g/cm$^3$. What percentage of formic acid molecules convert to formate anion:
\begin{qoptions}
\opt{A}{0.001 \%}
\opt{B}{0.002\%}
\opt{C}{0.003\%}
\opt{D}{0.004\%}
\end{qoptions}
\qn{3}
\qfig{p2_chem_s1_q03_1}
Decreasing order of shown angles will be
\begin{qoptions}
\opt{A}{$\theta_1>\theta_2>\theta_3$}
\opt{B}{$\theta_1>\theta_3>\theta_3$}
\opt{C}{$\theta_3>\theta_2>\theta_1$}
\opt{D}{$\theta_1=\theta_2=\theta_3$}
\end{qoptions}
\sect{SEC--2}
\qn{1} Between $\mathrm{ClO_2}$ and $\mathrm{Cl_2O}$
\begin{qoptions}
\opt{A}{The Cl $-$ O bond length would be larger in $\mathrm{Cl_2O}$}
\opt{B}{The bond angle would be larger in $\mathrm{Cl_2O}$}
\opt{C}{$\mathrm{ClO_2}$ is likely to be a coloured compound}
\opt{D}{$\mathrm{Cl_2O}$ is likely to be soluble in water}
\end{qoptions}
\qn{2} In the structure of $\beta,\beta'$-dichlorodiethyl sulphide, more popularly known as mustard gas:
\begin{qoptions}
\opt{A}{All C-atoms are $sp^3$ hybridised}
\opt{B}{S-atom is $sp^2$ hybridised}
\opt{C}{There is a total of 2 lone pair of electrons on one molecule}
\opt{D}{Order of bond strength :C--C $>$ C $-$ Cl $>$ C $-$ S}
\end{qoptions}
\qn{3} Which of the following carbonyl compounds(s) form(s) readily isolable hydrate.
\begin{qoptions}
\opt{A}{\qfigi{p2_chem_s2_q03_A}}
\opt{B}{$\mathrm{CCl_3-CHO}$}
\opt{C}{\qfigi{p2_chem_s2_q03_2}}
\opt{D}{\qfigi{p2_chem_s2_q03_D}}
\end{qoptions}
\qn{4} The correct statement(s) for the following reaction is/are.
\[\underset{\mathrm{I}}{\mathrm{[CoCl_3(NH_3)_3]}}+\mathrm{en}\longrightarrow\underset{\mathrm{II}}{\mathrm{[CoCl(en)(NH_3)_3]^{2+}}}+2\mathrm{Cl^-}\]
\begin{qoptions}
\opt{A}{If complex compound-I is a fac isomer, then complex ion-II will be fac-isomer w.r.t. $\mathrm{NH_3}$.}
\opt{B}{If the complex compound-I is mer isomer, then two isomers are produced.}
\opt{C}{If complex compound-I is a fac isomer, then complex ion-II exhibit optical activity.}
\opt{D}{If complex compound-I is a mer isomer then complex ion-II does not exhibit optical activity}
\end{qoptions}
\qn{5} Which of the given statement is correct :
\begin{qoptions}
\opt{A}{The C--Cl bond length in $\mathrm{CH_3CH_2Cl}$ is more than that in and CH=CH-Cl}
\opt{B}{The C $=$ C bond length in $\mathrm{CH_2=CH_2}$ is less than that in and $\mathrm{CH_2=CHCl}$}
\opt{C}{The C $-$ C bond length in $\mathrm{CH_2=CH_2}$, $\mathrm{CH_3CH_2-Cl}$ and $\mathrm{CH_2=CHCl}$ is equal}
\opt{D}{The dipole moment of $\mathrm{CH_3CH_2Cl}$ is more than that in $\mathrm{CH_2=CHCl}$}
\end{qoptions}
\sect{SEC--3}
\paraq{1}{2}
Redox reaction involves transfer of electrons between two chemical species. An unbalanced and incomplete example is shown below:\\
$\mathrm{N_2+H_3PO_2\longrightarrow HN_3+P_4O_6}$
\qn{1} If the above reaction were possible in either direction, maximum equivalent weight would belong to:
\begin{qoptions}
\opt{A}{$\mathrm{N_2}$}
\opt{B}{$\mathrm{H_3PO_2}$}
\opt{C}{$\mathrm{HN_3}$}
\opt{D}{$\mathrm{P_4O_6}$}
\end{qoptions}
\qn{2} If the above reaction were unidirectional towards right, what fraction of $\mathrm{P_4O_6}$ would be left unreacted if reaction were started with stoichiometric amount of reactants. (Assume complete reaction between $\mathrm{P_4O_6}$ and $\mathrm{H_2O}$, if formed in the above reaction)
\begin{qoptions}
\opt{A}{$\dfrac12$}
\opt{B}{$\dfrac34$}
\opt{C}{$\dfrac13$}
\opt{D}{$\dfrac23$}
\end{qoptions}
\sect{SEC--4}
\qn{1} Two inorganic compounds (A) and (B) are obtained by treating calcium orthophosphate with magnesium metal. (A) on hydrolysis produces a gas (C). The gas (C) reacts with $\mathrm{CuSO_4}$ (aq.) solution to give a black precipitate of compound (X). What is the molar mass of (X)?\\
{[}Atomic mass of P = 31 g, Ca = 40 g, Mg = 24 g, O = 16 g, Cu = 63.55].
\qn{2} In an ABCABC packing of Giant Buckies (C $-$ 240); Bucky babies - Chotu (C $-$ 32) and Motu (C $-$ 44) are doped in tetrahedral and octahedral voids respectively. In the unit cell thus created, all balls along any one body diagonal are removed and Bucky Onions (C $-$ 240 containing C $-$ 60 suspended inside it, with their centres coinciding) are fitted at both extremes. The number of C atoms finally present in the unit cell are.......?
\qn{3} Following graph is plotted between characteristics `X' and `Y' for three species A, B and C
\qfig{p2_chem_s4_q03_1}
Now consider the following statements, related to the above graph :\\
\textbf{Code-1.} If `X' is the velocity of electron and `Y' is $\dfrac1n$ (where `n' is $n$th Bohr orbit of three hydrogen like species A, B and C), then radii of these species must increase in the order C $<$ B $<$ A.\\
\textbf{Code-2}. If `X' is the de-Broglie wavelength ($\lambda$) of three charged species A, B and C (having same amount of charge) and `Y' is $=\dfrac1{\sqrt V}$ (where `V' is accelerating potential difference), then the mass of these species must increase in the order C $<$ B $<$ A.\\
\textbf{Code-3}. If `X' is the volume of same mole of three ideal gases A, B \& C and `Y' is the temperature in Kelvin, then the pressure of these gases must increase in the order C $<$ B $<$ A.\\
\textbf{Code-4}. If `X' is the pressure of the same mole of ideal gases A, B \& C and `Y' is $\dfrac1V$ (where V is the volume of these gases), then temperature of these gases must increase in the order : $T_C<T_B<T_A$. What is the correct increasing order of codes given , only for correct statements ? \textbf{(For exmaple, if only codes 1 \& 4 are correct, answer is 14.00)}
\sect{SEC--5}
\qn{1} The pH range of a basic indicator (InOH) is 3.4 to 4.6. For what minimum ratio of $\dfrac{[\mathrm{In^+}]}{[\mathrm{InOH}]}$ does the solution appear in the colour of $\mathrm{In^+}$. ($\log2=0.3$)
\qn{2}
\qfig[0.85]{p2_chem_s5_q02_1}
How many hydroxyl groups are present in X?
\qfullbegin
\qn{3} How many of the following compounds will give Cannizaro's reaction :
\qfig[0.78]{p2_chem_s5_q03_1}
\qfullend
\qn{4} An aromatic compound P($\mathrm{C_8H_{10}}$) is heated with alkaline $\mathrm{KMnO_4}$ to give a product which on heating with aqueous ammonia solution gives `Q'. Compound `Q' on heating with alkaline sodium hypochlorite gives a salt which on acidification gives product R ($\mathrm{C_7H_7O_2N}$). Compound `R' is treated with $\mathrm{NaNO_2}$ \& HCl and then with a base to produce a dipolar ion `S' which on heating produces an intermediate (T) which dimerises to produce the final product `U'. What is the degree of unsaturation of the compound (U)?
\qn{5} Number of alkenes of formula $\mathrm{C_4H_8}$ which give same product with HBr in presence as well as absence of peroxide.
\qn{6} 0.092 g of compound with molecular formula $\mathrm{C_3H_8O_3}$ on reaction with an excess of $\mathrm{CH_3MgI}$ gives 67 mL of methane at STP. What is the number of active hydrogen atoms present in a molecule of the compound?
\qn{7} The isomerisation of cyclopropane ($\mathrm{C_3H_6}$) to propene is believed to occur by the following mechanism:\\
Step-1 : $\mathrm{C_3H_6+C_3H_6\longrightarrow C_3H_6+C_3H_6^*}$\\
Step-2 : $\mathrm{C_3H_6^*\longrightarrow CH_3-CH=CH_2}$\\
At high pressure, step-1 is bidirectional and fast whereas step-2 is unidirectional and slow. What is the order of the reaction?
\qn{8} Ratio of frequency of revolution of electron in second excited state of $\mathrm{He^+}$ and second state of H is $\dfrac{2^x}{3^y}$. What is the value of `x $+$ y'?
\qfullbegin
\qn{9} How many of the following undergo solvolysis by 50\% aqueous ethanol, faster than
\qfig{p2_chem_s5_q09_1}
\qfullend
\qn{10} Consider the following (X):
\qfig{p2_chem_s5_q10_1}
If 0.5 moles of the compound (X) are hydrolysed in a alkaline medium and then heated. How many moles of $\mathrm{CO_2}$ are produced ?
\partsec{3}{MATHEMATICS}
\sect{SEC--1}
\qn{1} Solution of the differential equation $y^2\left(e^{x^3}+e^{y^3}\right)dy+x^2e^{x^3}\left(y^3-1\right)dx=0$ is\\
(Where $c$ is arbitrary constant)
\begin{qoptions}
\opt{A}{$e^{x^2}\left(y^3-1\right)+e^{y^3}=c$}
\opt{B}{$e^{x^3}\left(y^3+1\right)-e^{y^3}=c$}
\opt{C}{$e^{x^3}\left(y^3-1\right)+e^{y^3}=c$}
\opt{D}{$e^{x^3}\left(y^2-1\right)+e^{y^2}=c$}
\end{qoptions}
\qn{2} Number of positive integral solutions of $x_1x_2x_3x_4=3^{\alpha_1}5^{\beta_1}7^{\gamma_1}11^{\delta_1}$\\
(Where $\alpha_1,\beta_1,\gamma_1,\delta_1$ are integers such that $0\le\alpha_1\le1$, $0\le\beta_1\le1$, $0\le\gamma_1\le1$, $0\le\delta_1\le1$), is
\begin{qoptions}
\opt{A}{625}
\opt{B}{580}
\opt{C}{550}
\opt{D}{650}
\end{qoptions}
\qn{3} In $\Delta ABC$, $a$ is the arithmetic mean and $b$, $c$ are two geometric means of two positive distinct numbers, then the value of $4^4\left(\dfrac{\sin^{12}B+\sin^{12}C}{(\sin2A+\sin2B+\sin2C)^4}\right)$ can be equal to
\begin{qoptions}
\opt{A}{$\dfrac12$}
\opt{B}{$\dfrac14$}
\opt{C}{2}
\opt{D}{3}
\end{qoptions}
\sect{SEC--2}
\qn{1} Let $A_1,A_2,A_3,\ldots,A_n$ be $n$ points of a circle. Now, these points are joined to each other to form lines. So, find the number of new points of intersection of these lines if no two lines are parallel and no three lines are concurrent.
\begin{qoptions}
\opt{A}{$3.{}^nC_4$}
\opt{B}{${}^{{}^nC_2}C_2-{}^{n-1}C_2\,n$}
\opt{C}{$\dfrac{{}^{{}^nC_2}C_2-{}^{n-2}C_2}{2}$}
\opt{D}{${}^{{}^nC_2}C_2-n(n-1)$}
\end{qoptions}
\qn{2} Let $f(x)$ be a non-constant, twice differentiable function defined on $\mathbb R$ such that $y=f(x)$ is symmetric about line $x=1$ and $f(-1)=f'\!\left(\dfrac14\right)=f'\!\left(\dfrac12\right)=0$, then which of following is/are correct statement(s) ?
\begin{qoptions}
\opt{A}{$f''(x)=0$ has atleast four roots in the interval $(0,2)$}
\opt{B}{$\displaystyle\int_{-1}^{2}\left(x^5+x^3\right)f(1+x)\,dx<\int_0^2\dfrac1{1+2^{f(x)}}\,dx$}
\opt{C}{there exist at least one $C\in\left(\dfrac32,\dfrac74\right)$ such that $f'(C)+Cf''(C)=0$}
\opt{D}{for some $C\in\left(1,\dfrac32\right)$, $f'(C)=Cf''(C)$}
\end{qoptions}
\qn{3} Let $I=\displaystyle\int_0^{\pi}\dfrac{\cos\theta}{10-6\cos\theta}\,d\theta$, then which of following is/are correct statement(s)
\begin{qoptions}
\opt{A}{$I<1$}
\opt{B}{$I>0$}
\opt{C}{$I<0$}
\opt{D}{$I>1$}
\end{qoptions}
\qn{4} The line $L$ is drawn through the point $(1,1)$ to intersect $L_1: x+2y=1$ and $L_2: x+2y=3$ at $P$ and $Q$ respectively. Line perpendicular to $L$ from $P$ intersects $L_2$ at $R$ and line perpendicular to $L$ from $Q$ intersects $L_1$ and $S$. If the area of parallelogram $PQRS$ is least, then which of the following is/are correct ?
\begin{qoptions}
\opt{A}{the equation of line $L$ can be $y+3x=4$}
\opt{B}{the equation of line $L$ can be $3y-x=2$}
\opt{C}{the least area of $PQRS$ is $\dfrac85$ square unit}
\opt{D}{the equation of line $L$ can be $3y-x=4$}
\end{qoptions}
\qn{5} Let $ABCD$ be a rectangle inscribed in a circle whose one diameter is along line $y=x$. If coordinates of $A$ and $B$ are $(0,2)$ and $(4,3)$ respectively, then which of the following is/are correct ?
\begin{qoptions}
\opt{A}{the coordinates of point $D$ are $\left(\dfrac15,\dfrac65\right)$}
\opt{B}{the coordinates of point $D$ are $(4,6)$}
\opt{C}{the coordinates of $C$ are $\left(\dfrac{21}5,\dfrac{11}5\right)$}
\opt{D}{the coordinates of $C$ are $(6,4)$}
\end{qoptions}
\sect{SEC--3}
\paraq{1}{2}
Five persons shoot each other such that no one shoots himself, every one shoots exactly once and no shot is wasted, then
\qn{1} The number of ways such that all persons are hit is
\begin{qoptions}
\opt{A}{60}
\opt{B}{44}
\opt{C}{120}
\opt{D}{36}
\end{qoptions}
\qn{2} The number of ways in which exactly one person is not hit is
\begin{qoptions}
\opt{A}{420}
\opt{B}{180}
\opt{C}{210}
\opt{D}{540}
\end{qoptions}
\sect{SEC--4}
\qn{1} $P_1,P_2,P_3,\ldots..,P_8$ are 8 players of equal strength who play with each other in a knockout tournament. The probability that $P_1$ wins against $P_2$ and $P_3$ to reach the final and also wins the tournament is $p$, then $42p$ is equal to
\qn{2} Let $\displaystyle\lim_{n\to\infty}\sum_{k=1}^{n}\left|e^{i\frac{2k\pi}n}-e^{i\frac{(k-1)2\pi}n}\right|=L$, then $\dfrac L{10\pi}$ is equal to
\qn{3} The arithmetic mean of all integral values of real parameter $\alpha$ for which $|z-(\alpha^2-7\alpha+13+i)|=1$ and $\arg z\geqslant\dfrac\pi2$ is satisfied for at least one $z$ is equal to
\sect{SEC--5}
\qn{1} If the number of onto functions from $\{1,2,3,4\}\to\{2,3,4\}$ such that $f(i)\ne i$ for $i=1,2,3,4$ is $k$, then $\dfrac k4$ is equal to
\qn{2} Let the area enclosed by the curve $x^4+y^4=2xy$ be $A$ square units, then the value of $[A]$ is equal to (where $[.]$ denotes the greatest integer function)
\qn{3} Let the probability of arranging 8 boys ($B_1,B_2,\ldots..,B_8$) and 5 girls linearly such that $B_1$ and $B_2$ are not together and exactly four girls are together be $p$, then $[99p]$ is equal to (where $[.]$ denotes the greatest integer function)
\qn{4} If $3\sin\alpha\sin\beta+4\cos\beta+5\cos\alpha\sin\beta=5\sqrt2\quad\forall\ \alpha,\beta\in\mathbb R$, then $|\operatorname{adj}\operatorname{adj}A|$ is equal to $\left\{\text{Where }A=\begin{bmatrix}\sin\alpha&-\frac12\\[2pt]\frac12&\tan\beta\end{bmatrix}\right\}$
\qn{5} In $\Delta ABC$ if $r_1=3r$, then the value of $9\left(\tan\dfrac A2\tan\dfrac B2+\tan\dfrac A2\tan\dfrac C2\right)$ is equal to (where $r_1$ and $r$ are the ex-radius opposite $A$ and the inradius of $\Delta ABC$ respectively)
\qn{6} If $\alpha,\beta,\gamma$ are the roots of $x^3+x^2-2x+1=0$, then the value of $\sqrt\Delta$ is equal to
\[\text{where }\Delta=\begin{vmatrix}\gamma^2&\beta^2&2\beta\gamma-\alpha^2\\2\alpha\gamma-\beta^2&\alpha^2&\gamma^2\\\alpha^2&2\alpha\beta-\gamma^2&\beta^2\end{vmatrix}\]
\qn{7} The value of $\dfrac{\displaystyle\sum_{r=1}^{k}2^r\,{}^nC_r\,{}^{k-1}C_{r-1}}{\displaystyle\sum_{r=0}^{k}{}^nC_r\,{}^{n+k-r-1}C_{n-1}}$ (where $n\ge k$) is equal to
\qn{8} Let in $\Delta ABC$, $\angle A$ is $\dfrac\pi2$ and $D$ is point on $BC$ such that $AD$ is altitude. If $E$, $F$ and $I$ are incentre of $\Delta ABD$, $\Delta ADC$ and $\Delta ABC$ respectively, then the value of $\dfrac{AI}{EF}$ is equal to
\qn{9} The value of $m^m-\displaystyle\sum_{r=1}^{m}\dfrac{(m+1)(r-1)m^{r-1}}{r\,{}^mC_r}$ is equal to
\qn{10} A variable chord of the hyperbola $\dfrac{x^2}4-\dfrac{y^2}8=1$, subtends a right angle at the centre of the hyperbola. If this chord touches a fixed circle concentric with the hyperbola, then the square of radius of the fixed circle is equal to

\end{multicols}
\end{document}
